% GENERATED FILE. Edit canonical lessons, then run npm run generate:books.
% canonical-source-hash: fnv1a64:b70f02351970ba5b
% canonical-lessons: GU-C397-dhamal, GU-C397-umbro, GU-C397-patlo, GU-C397-hadselvu, GU-C397-bharavu

\chapter{The Rush at the Door}
\label{ch:gu-the-rush-at-the-door}
\hlchaptermodality{397}

\begin{chapteropening}
\textbf{By the end of this chapter:} \emph{I can get through a crowd at the carriage door.}

Complete the last lesson of chapter 397: I can get through a crowd at the carriage door.
\end{chapteropening}

\section[\texorpdfstring{\gu{ધમાલ} (dhamāl)}{ધમાલ (dhamāl)}]{\gu{ધમાલ} (dhamāl) — a commotion; a bustle}

\label{lesson:GU-C397-dhamal}



\begin{quote}
Before the new one: say the Gujarati for to ring, then the Gujarati for a large bell.
\end{quote}

\subsection*{You'll want to know: \gu{ધમાલ}}
\textbf{\gu{ધમાલ}} — \emph{dhamāl} — \textquotedblleft{}a commotion; a bustle\textquotedblright{}.

\begin{quote}
\gu{ટ્રેન} \gu{આવી} \gu{ત્યારે} \gu{પ્લેટફોર્મ} \gu{પર} \gu{ધમાલ} \gu{થઈ}.
\end{quote}

\begin{quote}
\emph{ṭren āvī tyāre pleṭphorm par dhamāl thaī.}
\end{quote}

\begin{quote}
When the train came in there was a commotion on the platform.
\end{quote}

A feminine noun.

\begin{cousinweb}
\textbf{\gu{ધમાલ}} (\emph{dhamāl}) No dictionary we use records an older form, so its history is not settled.
\end{cousinweb}

\begin{grammarlens}[title={What you've built}]
One more word for getting around.
\end{grammarlens}

\subsection*{Your turn}
\begin{itemize}
\raggedright
  \item \emph{Say it:} \emph{dhamāl}
  \item \emph{Say it:} \emph{dhamāl}, once more
  \item \emph{Say it:} \emph{ṭren āvī tyāre pleṭphorm par dhamāl thaī.} — When the train came in there was a commotion on the platform.
  \item \emph{Recall:} say the Gujarati for to ring, then the Gujarati for a large bell, then say \emph{dhamāl} again
\end{itemize}

\begin{tcolorbox}[breakable,colback=teal!4,colframe=teal!35!black,title={Before you move on}]
What does \gu{ધમાલ} mean? (\textquotedblleft{}A commotion; a bustle\textquotedblright{}.) Say it once more.
\end{tcolorbox}
\section[\texorpdfstring{\gu{ઉંબરો} (umbro)}{ઉંબરો (umbro)}]{\gu{ઉંબરો} (umbro) — a threshold (of a door)}

\label{lesson:GU-C397-umbro}



\begin{quote}
Before the new one: say the Gujarati for a large bell, then the Gujarati for a commotion.
\end{quote}

\subsection*{You'll want to know: \gu{ઉંબરો}}
\textbf{\gu{ઉંબરો}} — \emph{umbro} — \textquotedblleft{}a threshold (of a door)\textquotedblright{}.

\begin{quote}
\gu{ઘરનો} \gu{ઉંબરો} \gu{ઓળંગીને} \gu{તે} \gu{બહાર} \gu{ગયો}.
\end{quote}

\begin{quote}
\emph{gharno umbro oḷangīne te bahār gayo.}
\end{quote}

\begin{quote}
He stepped over the threshold of the house and went out.
\end{quote}

A masculine noun.

\begin{cousinweb}
\textbf{\gu{ઉંબરો}} (\emph{umbro}) is a tadbhava of Sanskrit \emph{umbara}, the timber across the top of a door frame; CDIAL lists the Gujarati word, \textquotedblleft{}a threshold, a sill\textquotedblright{} (CDIAL 2341): the beam has moved from the top of the door to the bottom.
\end{cousinweb}

\begin{grammarlens}[title={What you've built}]
One more word for getting around.
\end{grammarlens}

\subsection*{Your turn}
\begin{itemize}
\raggedright
  \item \emph{Say it:} \emph{umbro}
  \item \emph{Say it:} \emph{umbro}, once more
  \item \emph{Say it:} \emph{gharno umbro oḷangīne te bahār gayo.} — He stepped over the threshold of the house and went out.
  \item \emph{Recall:} say the Gujarati for a large bell, then the Gujarati for a commotion, then say \emph{umbro} again
\end{itemize}

\begin{tcolorbox}[breakable,colback=teal!4,colframe=teal!35!black,title={Before you move on}]
What does \gu{ઉંબરો} mean? (\textquotedblleft{}A threshold (of a door)\textquotedblright{}.) Say it once more.
\end{tcolorbox}
\section[\texorpdfstring{\gu{પાટલો} (pāṭlo)}{પાટલો (pāṭlo)}]{\gu{પાટલો} (pāṭlo) — a low wooden seat; a bench}

\label{lesson:GU-C397-patlo}



\begin{quote}
Before the new one: say the Gujarati for a commotion, then the Gujarati for a threshold (of a door).
\end{quote}

\subsection*{You'll want to know: \gu{પાટલો}}
\textbf{\gu{પાટલો}} — \emph{pāṭlo} — \textquotedblleft{}a low wooden seat; a bench\textquotedblright{}.

\begin{quote}
\gu{દાદી} \gu{પાટલા} \gu{પર} \gu{બેસીને} \gu{જમે} \gu{છે}.
\end{quote}

\begin{quote}
\emph{dādī pāṭlā par besīne jame che.}
\end{quote}

\begin{quote}
Grandmother sits on a low wooden seat to eat.
\end{quote}

A masculine noun.

\begin{cousinweb}
\textbf{\gu{પાટલો}} (\emph{pāṭlo}) is a tadbhava of Sanskrit \emph{paṭṭa}, \textquotedblleft{}a slab, a board\textquotedblright{}; CDIAL lists the Gujarati word, \textquotedblleft{}a bench\textquotedblright{} (CDIAL 7699). \gu{પાટિયું}, a plank, comes from the same word.
\end{cousinweb}

\begin{grammarlens}[title={What you've built}]
One more word for getting around.
\end{grammarlens}

\subsection*{Your turn}
\begin{itemize}
\raggedright
  \item \emph{Say it:} \emph{pāṭlo}
  \item \emph{Say it:} \emph{pāṭlo}, once more
  \item \emph{Say it:} \emph{dādī pāṭlā par besīne jame che.} — Grandmother sits on a low wooden seat to eat.
  \item \emph{Recall:} say the Gujarati for a commotion, then the Gujarati for a threshold (of a door), then say \emph{pāṭlo} again
\end{itemize}

\begin{tcolorbox}[breakable,colback=teal!4,colframe=teal!35!black,title={Before you move on}]
What does \gu{પાટલો} mean? (\textquotedblleft{}A low wooden seat; a bench\textquotedblright{}.) Say it once more.
\end{tcolorbox}
\section[\texorpdfstring{\gu{હડસેલવું} (haḍselvũ)}{હડસેલવું (haḍselvũ)}]{\gu{હડસેલવું} (haḍselvũ) — to shove; to elbow aside}

\label{lesson:GU-C397-hadselvu}



\begin{quote}
Before the new one: say the Gujarati for a threshold (of a door), then the Gujarati for a low wooden seat.
\end{quote}

\subsection*{You'll want to know: \gu{હડસેલવું}}
\textbf{\gu{હડસેલવું}} — \emph{haḍselvũ} — \textquotedblleft{}to shove; to elbow aside\textquotedblright{}.

\begin{quote}
\gu{ભીડમાં} \gu{બધા} \gu{માણસો} \gu{બીજાને} \gu{હડસેલતા} \gu{હતા}.
\end{quote}

\begin{quote}
\emph{bhīḍmā̃ badhā māṇso bījāne haḍseltā hatā.}
\end{quote}

\begin{quote}
In the crowd everyone was shoving everyone else.
\end{quote}

Past, with \gu{તેણે} (he or she …): \gu{તેણે} \gu{થેલી} \gu{હડસેલી}, \emph{teṇe thelī haḍselī}.

\begin{cousinweb}
\textbf{\gu{હડસેલવું}} (\emph{haḍselvũ}) goes back to \emph{*haṭ}, \textquotedblleft{}to move or cry out violently\textquotedblright{}, a form CDIAL reconstructs; CDIAL lists the Gujarati verb, \textquotedblleft{}to push violently\textquotedblright{} (CDIAL 13942).
\end{cousinweb}

\begin{grammarlens}[title={What you've built}]
One more word for getting around.
\end{grammarlens}

\subsection*{Your turn}
\begin{itemize}
\raggedright
  \item \emph{Say it:} \emph{haḍselvũ}
  \item \emph{Say it:} \emph{haḍselvũ}, once more
  \item \emph{Say it:} \emph{bhīḍmā̃ badhā māṇso bījāne haḍseltā hatā.} — In the crowd everyone was shoving everyone else.
  \item \emph{Recall:} say the Gujarati for a threshold (of a door), then the Gujarati for a low wooden seat, then say \emph{haḍselvũ} again
\end{itemize}

\begin{tcolorbox}[breakable,colback=teal!4,colframe=teal!35!black,title={Before you move on}]
What does \gu{હડસેલવું} mean? (\textquotedblleft{}To shove; to elbow aside\textquotedblright{}.) Say it once more.
\end{tcolorbox}
\section[\texorpdfstring{\gu{ભરાવું} (bharāvũ)}{ભરાવું (bharāvũ)}]{\gu{ભરાવું} (bharāvũ) — to fill up; to become full}

\label{lesson:GU-C397-bharavu}



\begin{quote}
Before the new one: say the Gujarati for a low wooden seat, then the Gujarati for to shove.
\end{quote}

\subsection*{You'll want to know: \gu{ભરાવું}}
\textbf{\gu{ભરાવું}} — \emph{bharāvũ} — \textquotedblleft{}to fill up; to become full\textquotedblright{}.

\begin{quote}
\gu{સાંજે} \gu{આ} \gu{ટ્રેન} \gu{તરત} \gu{ભરાઈ} \gu{જાય} \gu{છે}.
\end{quote}

\begin{quote}
\emph{sā̃je ā ṭren tarat bharāī jāy che.}
\end{quote}

\begin{quote}
In the evening this train fills up at once.
\end{quote}

Past (it …): \gu{તે} \gu{ડબ્બો} \gu{ભરાઈ} \gu{ગયો}, \emph{te ḍabbo bharāī gayo}.

\begin{cousinweb}
\textbf{\gu{ભરાવું}} (\emph{bharāvũ}) is the intransitive partner of \gu{ભરવું}, \textquotedblleft{}to fill\textquotedblright{}, a tadbhava of Sanskrit \emph{bhárati}, \textquotedblleft{}bears, fills\textquotedblright{} (CDIAL 9397): what you fill, \gu{ભરવું}; what fills up by itself, \gu{ભરાવું}.
\end{cousinweb}

\begin{grammarlens}[title={What you've built}]
One more word for getting around.
\end{grammarlens}

\subsection*{Your turn}
\begin{itemize}
\raggedright
  \item \emph{Say it:} \emph{bharāvũ}
  \item \emph{Say it:} \emph{bharāvũ}, once more
  \item \emph{Say it:} \emph{sā̃je ā ṭren tarat bharāī jāy che.} — In the evening this train fills up at once.
  \item \emph{Recall:} say the Gujarati for a low wooden seat, then the Gujarati for to shove, then say \emph{bharāvũ} again
\end{itemize}

\begin{tcolorbox}[breakable,colback=teal!4,colframe=teal!35!black,title={Before you move on}]
What does \gu{ભરાવું} mean? (\textquotedblleft{}To fill up; to become full\textquotedblright{}.) Say it once more.
\end{tcolorbox}
